\begin{insight}{Jordan链解释核为何逐级增长}
在一条长度$s$的幂零链上，一次作用把核端的首向量$v_1$送到零，$k$次作用能消去$\min(k,s)$个独立方向。把各条链的贡献相加，就得到核维数。增长量不增，是因为“长度至少$k+1$”的链集合包含于“长度至少$k$”的链集合；这不表示恰为某长度的块数必须单调。
\end{insight}
\begin{figure}[H]\centering
\begin{tikzpicture}[scale=1.05]
\foreach \xx/\txt in {0/v_3,1.2/v_2,2.4/v_1}{\node[draw=Teal,fill=Teal!6,minimum width=8mm]at(\xx,2){$\txt$};}
\foreach \xx in {0.4,1.6,2.8}{\draw[-{Stealth},Teal](\xx,2)--(\xx+.4,2);}\node at(3.6,2){$0$};
\node[draw=Gold,fill=Gold!6,minimum width=8mm]at(1.2,1){$w_2$};\node[draw=Gold,fill=Gold!6,minimum width=8mm]at(2.4,1){$w_1$};
\draw[-{Stealth},Gold](1.6,1)--(2,1);\draw[-{Stealth},Gold](2.8,1)--(3.2,1);\node at(3.6,1){$0$};
\node[draw=Ink,fill=Ink!4,minimum width=8mm]at(2.4,0){$z_1$};\draw[-{Stealth},Ink](2.8,0)--(3.2,0);\node at(3.6,0){$0$};
\node[lanote]at(1.8,-.75){三条链的长度为$3,2,1$};
\draw[laaxis](5,-.05)--(8.5,-.05);
\foreach \xx/\hh/\num/\lab in {5.5/.9/3/1,6.6/1.5/5/2,7.7/1.8/6/3}{\filldraw[fill=Teal!15,draw=Teal](\xx-.3,0)rectangle(\xx+.3,\hh);\node[above]at(\xx,\hh){$\num$};\node[font=\small,below]at(\xx,-.12){$d_{\lab}$};}
\node[lanote]at(6.6,-.75){核维数：$3,5,6$};
\end{tikzpicture}
\caption{每一条箭头表示作用一次$N$；核增长量为$3,2,1$，依次是长度至少1、2、3的链数。}\label{fig:jordan-kernel-growth}
\end{figure}
\begin{clevertrick}{看核的差分，不猜Jordan形}
给出$d_0=0,d_1=3,d_2=5,d_3=6$且此后稳定，令$b_k=d_k-d_{k-1}$，逐项得$b_1=3,b_2=2,b_3=1,b_4=0$。一条长度$s$的链对$d_k$贡献$\min(k,s)$，所以对$b_k$贡献1恰在$s\ge k$时成立；故$b_k$数的是长度至少$k$的块。恰为长度$k$的块数则是$b_k-b_{k+1}$，本例为
\[b_1-b_2=1,\qquad b_2-b_3=1,\qquad b_3-b_4=1.\]
所以块为$J_1(0),J_2(0),J_3(0)$各一个。总大小$1+2+3=6$；反算核维数为$1+1+1=3,1+2+2=5,1+2+3=6$，验证全部数据，无需枚举六阶矩阵的所有分拆。
\textbf{易错提示：}核维数必须不减，增长量必须不增；若计算得$d_1=1,d_2=3$，就与链计数矛盾，应返回检查消元。
\end{clevertrick}
\begin{insight}{SVD是“换正交方向、分别缩放、再换方向”}
$A=U\Sigma V\T$中，$V\T$先转到输入主方向，$\Sigma$按奇异值缩放，$U$再把输出主方向摆回去。对二维可逆矩阵，单位圆像是非退化椭圆，其最长、最短半轴为两个正奇异值；若矩阵奇异，图形会退化为线段或一点，不能仍按两个正半轴读取。非奇异方阵的条件数比较最大与最小缩放率，而不是面积。
\end{insight}
\begin{figure}[H]\centering
\begin{tikzpicture}[scale=1.15]
\begin{scope}
\draw[laaxis](-1.4,0)--(1.5,0);\draw[laaxis](0,-1.4)--(0,1.5);\draw[Teal,thick](0,0)circle(1);
\node[lanote]at(0,-1.85){输入单位圆};
\end{scope}
\draw[-{Stealth},Gold,thick](1.8,0)--(3.2,0)node[midway,above]{$\diag(3,1)$};
\begin{scope}[xshift=6.8cm]
\draw[laaxis](-3.4,0)--(3.5,0);\draw[laaxis](0,-1.4)--(0,1.5);\draw[Teal,thick](0,0)ellipse(3 and 1);
\draw[lavec,Ink](0,0)--(3,0)node[midway,below]{$\sigma_1=3$};\draw[lavec,Gold](0,0)--(0,1)node[above left]{$\sigma_2=1$};
\node[lanote]at(0,-1.85){输出椭圆，$\kappa_2=3$};
\end{scope}
\end{tikzpicture}
\caption{奇异值是单位球各主方向的缩放率；一般矩阵还会在两侧作正交方向变换。}\label{fig:svd-ellipse}
\end{figure}
\begin{clevertrick}{用尺度反例区分行列式和条件数}
先取$0<\varepsilon<1$，$A=\varepsilon I_2$。有$\det A=\varepsilon^2$、$A\T A=\varepsilon^2I_2$，所以两个奇异值都为$\varepsilon$，$\kappa_2(A)=\frac{\varepsilon}{\varepsilon}=1$。当$\varepsilon\to0^+$，行列式趋于零，但在每个非零$\varepsilon$下条件数都为1。
再取$M\ge1$，$B=\diag(M,\frac{1}{M})$。有$\det B=M\cdot(\frac{1}{M})=1$，$B\T B=\diag(M^2,M^{-2})$，所以最大、最小奇异值为$M,\frac{1}{M}$，
\[\kappa_2(B)=\dfrac{M}{\frac{1}{M}}=M^2\longrightarrow\infty\quad(M\to\infty).\]
前者在各方向均匀缩小，后者沿一轴拉长、另一轴压缩；面积分别记录奇异值的乘积，而条件数比较最大与最小的比值。图中$\diag(3,1)$同理有奇异值3、1，故$\kappa_2=3$，不能仅凭面积判断方向伸缩是否均衡。
\end{clevertrick}
